
\documentclass{article}
\usepackage{colortbl}
\usepackage{makecell}
\usepackage{multirow}
\usepackage{supertabular}

\begin{document}

\newcounter{utterance}

\centering \large Interaction Transcript for game `dond', experiment `semi\_en', episode 1 with Qwen3.5{-}9B.
\vspace{24pt}

{ \footnotesize  \setcounter{utterance}{1}
\setlength{\tabcolsep}{0pt}
\begin{supertabular}{c@{$\;$}|p{.15\linewidth}@{}p{.15\linewidth}p{.15\linewidth}p{.15\linewidth}p{.15\linewidth}p{.15\linewidth}}
    \# & $\;$A & \multicolumn{4}{c}{Game Master} & $\;\:$B\\
    \hline

    \theutterance \stepcounter{utterance}  
    & & \multicolumn{4}{p{0.6\linewidth}}{
        \cellcolor[rgb]{0.9,0.9,0.9}{
            \makecell[{{p{\linewidth}}}]{
                \texttt{\tiny{[P1$\langle$GM]}}
                \texttt{You are playing a negotiation game in which you have to agree on how to divide a set of items among you and another player.} \\
\\ 
\texttt{Rules:} \\
\texttt{(a) You and the other player are given a set of items. Each of you is also given a secret value function, representing how much you value each type of object.} \\
\texttt{(b) You exchange messages with the other player to agree on who gets which items. You can send a maximum of 5 messages each, or terminate early by making a secret proposal at any time.} \\
\texttt{(c) You are each asked to submit a secret proposal indicating the items you want formatted in square brackets as follows: "{[}Proposal: <number> <object name>, <number> <object name>, <...>{]}"} \\
\texttt{(d) If your proposals are complementary, i.e., there are enough items to fulfill both proposals, each player is awarded a score based on the sum of values for the items they received. Otherwise, both of you get zero points.} \\
\texttt{(e) Your goal is to maximize the score you receive.} \\
\\ 
\texttt{Let us start.} \\
\\ 
\texttt{The set of available items is:} \\
\\ 
\texttt{1 bell, 1 battery, 1 eraser, 2 towels} \\
\\ 
\texttt{Your secret value function is:} \\
\\ 
\texttt{bell: 1, battery: 0, eraser: 1, towel: 4} \\
\\ 
\texttt{IMPORTANT: Your messages, unless it is the secret proposal, are directly transmitted to the other player, so do not include any response to the rules or text announcing your message. To make a secret proposal, use the indicated format. Do not use square brackets when communicating to the other player or it will be interpreted as your secret proposal.} \\
            }
        }
    }
    & \\ \\

    \theutterance \stepcounter{utterance}  
    & \multicolumn{4}{p{0.6\linewidth}}{
        \cellcolor[rgb]{0.9,0.9,0.9}{
            \makecell[{{p{\linewidth}}}]{
                \texttt{\tiny{[P1$\rangle$GM]}}
                \texttt{I value the towels the most (4 points each), followed by the bell and eraser (1 point each). The battery is worthless to me (0 points). Since there are only 2 towels available, I need to ensure I get both. To maximize my chances of agreement, I should propose taking the two towels and offer the other player the items they might value (the bell and eraser), while leaving the battery for them if they want it, or just giving it to them as a "free" item to make the deal more attractive.} \\
\\ 
\texttt{However, I don't know the other player's values. A safe strategy is to propose a bundle that includes the high{-}value items for me (the towels) and offers the remaining high{-}value items (bell, eraser) to the other player, effectively splitting the non{-}battery items. If they value the bell or eraser, they will accept. If they value the towels more, they might counter{-}offer, but I can only make 5 messages. Let's start with a clear division proposal.} \\
\\ 
\texttt{I will propose:} \\
\texttt{{-} Me: 2 towels} \\
\texttt{{-} You: 1 bell, 1 eraser} \\
\\ 
\texttt{This leaves the battery unassigned in my proposal, which might be an issue if the system requires all items to be assigned or if "complementary" means the union of proposals must equal the set. The rules say "If your proposals are complementary, i.e., there are enough items to fulfill both proposals". This implies the sum of items requested by both must not exceed the available supply. It doesn't explicitly say every item must be taken, but usually, in these games, leaving items unassigned is fine as long as no one asks for more than exists.} \\
\\ 
\texttt{Wait, looking at the format: `{[}Proposal: <number> <object name>, ...{]}`.} \\
\texttt{My proposal should list what *I* want.} \\
\texttt{If I say: `{[}Proposal: 2 towel{]}`} \\
\texttt{And you say: `{[}Proposal: 1 bell, 1 eraser{]}`} \\
\texttt{Total requested: 2 towels, 1 bell, 1 eraser. Available: 2 towels, 1 bell, 1 eraser, 1 battery.} \\
\texttt{This is complementary. I get 2 towels (Score: 8). You get bell + eraser.} \\
\\ 
\texttt{Let's try to be even more generous to ensure acceptance. I'll offer you the bell and eraser, and I'll take the towels. I won't mention the battery in my request, implying I don't want it (since it's worth 0 to me).} \\
\\ 
\texttt{Message 1:} \\
\texttt{I want the two towels. I offer you the bell and the eraser.} \\
\\ 
\texttt{Let's refine the message to be direct and clear, avoiding any ambiguity about the battery.} \\
\\ 
\texttt{Message:} \\
\texttt{I propose taking the 2 towels. I offer you the 1 bell and the 1 eraser.} \\
            }
        }
    }
    & & \\ \\

    \theutterance \stepcounter{utterance}  
    & & & \multicolumn{2}{p{0.3\linewidth}}{
        \cellcolor[rgb]{0.9,0.9,0.9}{
            \makecell[{{p{\linewidth}}}]{
                \texttt{\tiny{[GM$|$GM]}}
                \texttt{proposal does not include number/name pairs} \\
            }
        }
    }
    & & \\ \\

\end{supertabular}
}

\end{document}
